% ============================================================
\chapter*{Preface}
\addcontentsline{toc}{chapter}{Preface}
% ============================================================

Partial differential equations (PDEs) form one of the fundamental pillars of
applied mathematics and mathematical physics. Since the foundational work of
Euler, d'Alembert, Fourier, and Laplace in the 18th century, they have played
a central role in modelling natural phenomena: heat diffusion, wave propagation,
fluid flow, material elasticity, electromagnetism, and more recently in finance,
biology, and image processing.

\section*{Course objectives}

This course is aimed at third-year undergraduate students (L3 level) in
mathematics. It assumes the following prerequisites:
\begin{itemize}
    \item Real analysis: Lebesgue integration, $L^p$ spaces, modes of convergence.
    \item Elementary functional analysis: Banach spaces, Hilbert spaces.
    \item Ordinary differential equations (ODEs): Cauchy--Lipschitz theorem.
    \item Linear algebra: diagonalization, bilinear forms.
\end{itemize}

The main objective is to provide a rigorous yet accessible introduction to
classical and modern methods for solving PDEs. More precisely, upon
completion of this course, the student will be able to:
\begin{enumerate}
    \item Classify second-order PDEs (elliptic, parabolic, hyperbolic) and
          understand the physical implications of each type.
    \item Solve classical PDEs explicitly using the method of characteristics,
          separation of variables, and Fourier series.
    \item State and prove fundamental qualitative properties: maximum principle,
          regularity, asymptotic behaviour.
    \item Use the Fourier transform to solve PDEs on $\R^n$.
    \item Understand the modern functional framework: Sobolev spaces and
          variational formulations.
\end{enumerate}

\section*{Course organization}

The course is organized into ten chapters, following a natural pedagogical
progression:

\begin{description}
    \item[Chapter 1:] \textbf{Introduction and classification.} Basic vocabulary
        (order, linearity, boundary conditions) and classification of
        second-order PDEs.

    \item[Chapter 2:] \textbf{Transport equation.} The first PDE studied in
        detail, introducing the method of characteristics in a simple setting.

    \item[Chapter 3:] \textbf{Separation of variables and Fourier series.}
        Development of the central tool of separation of variables, supported
        by the theory of Fourier series.

    \item[Chapter 4:] \textbf{Heat equation.} Complete study: fundamental
        solution, maximum principle, regularity, asymptotic behaviour.

    \item[Chapter 5:] \textbf{Wave equation.} d'Alembert's formula, Huygens'
        principle, energy and finite speed of propagation.

    \item[Chapter 6:] \textbf{Laplace equation and harmonic functions.}
        Mean value properties, regularity, Harnack inequality.

    \item[Chapter 7:] \textbf{Maximum principle.} Weak and strong versions for
        elliptic and parabolic operators.

    \item[Chapter 8:] \textbf{Fourier transform and applications.}
        The Fourier transform as a tool for solving PDEs on $\R^n$.

    \item[Chapter 9:] \textbf{Introduction to Sobolev spaces.}
        Weak derivatives, $H^1$ and $H^1_0$ spaces, Sobolev and Poincar\'e
        inequalities.

    \item[Chapter 10:] \textbf{Variational formulation.}
        Lax--Milgram theorem, elliptic problems, introduction to finite
        elements.
\end{description}

\section*{Conventions and notation}

Throughout this course, we use the following notation:
\begin{itemize}
    \item $\Omega$ denotes a bounded open subset of $\R^n$ with boundary
          $\partial\Omega$ sufficiently regular (at least piecewise $C^1$,
          unless otherwise stated).
    \item $u_t = \frac{\partial u}{\partial t}$, $u_x = \frac{\partial u}{\partial x}$,
          $u_{xx} = \frac{\partial^2 u}{\partial x^2}$.
    \item $\Delta u = \sum_{i=1}^{n} \frac{\partial^2 u}{\partial x_i^2}$ is
          the Laplacian of $u$.
    \item $\nabla u = \left(\frac{\partial u}{\partial x_1}, \ldots,
          \frac{\partial u}{\partial x_n}\right)$ is the gradient.
    \item $Du$ denotes the Jacobian matrix, $D^2 u$ the Hessian matrix.
    \item $C^k(\Omega)$, $C^\infty(\Omega)$, $C_c^\infty(\Omega)$: spaces of
          $C^k$, smooth, and compactly supported smooth functions.
    \item $L^p(\Omega)$: Lebesgue space, $1 \leq p \leq \infty$.
    \item $H^s(\Omega) = W^{s,2}(\Omega)$: Sobolev space.
    \item $\inner{\cdot}{\cdot}$ denotes the inner product in $L^2$ or in a
          Hilbert space.
    \item $\norm{\cdot}$ denotes a norm (context specifies which one).
    \item Multi-indices: for $\alpha = (\alpha_1, \ldots, \alpha_n) \in \N^n$,
          $\abs{\alpha} = \alpha_1 + \cdots + \alpha_n$ and
          $D^\alpha u = \frac{\partial^{\abs{\alpha}} u}{\partial x_1^{\alpha_1}
          \cdots \partial x_n^{\alpha_n}}$.
\end{itemize}

\section*{How to use this course}

Each chapter contains:
\begin{itemize}
    \item \textbf{Definitions} and \textbf{theorems} with detailed proofs.
    \item \textbf{Examples} worked out completely to illustrate the methods.
    \item \textbf{Intuition} boxes to develop geometric and physical understanding.
    \item \textbf{Warning} boxes to highlight common pitfalls.
    \item \textbf{Exercises} of increasing difficulty, some with hints.
\end{itemize}

Proofs marked with an asterisk (*) are beyond the syllabus and are offered
for students wishing to deepen their understanding.

\section*{Physical motivation}

PDEs arise naturally whenever one models a phenomenon depending on several
independent variables. Let us consider some fundamental examples.

\paragraph{Heat diffusion.}
If $u(x,t)$ represents the temperature at point $x$ at time $t$ in a
conducting body, Fourier's law states that the heat flux is proportional to
the temperature gradient: $\vec{q} = -k\nabla u$. Combined with energy
conservation, one obtains the heat equation:
\[
    \frac{\partial u}{\partial t} = k \Delta u.
\]

\paragraph{Wave propagation.}
Vibrations of a string, a membrane, or acoustic waves are described by the
wave equation:
\[
    \frac{\partial^2 u}{\partial t^2} = c^2 \Delta u,
\]
where $c$ is the propagation speed.

\paragraph{Electrostatics.}
The electric potential $u$ in a charge-free region satisfies Laplace's
equation:
\[
    \Delta u = 0.
\]
In the presence of a charge density $\rho$, one obtains Poisson's equation
$\Delta u = -\rho/\varepsilon_0$.

\paragraph{Fluid mechanics.}
The Navier--Stokes equations, describing the motion of a viscous
incompressible fluid, form a system of nonlinear PDEs:
\[
    \frac{\partial \vec{v}}{\partial t} + (\vec{v} \cdot \nabla)\vec{v}
    = -\frac{1}{\rho}\nabla p + \nu \Delta \vec{v}, \qquad
    \nabla \cdot \vec{v} = 0.
\]
The existence and regularity of solutions in three dimensions remains one of
the Millennium Prize Problems.

\paragraph{Quantum mechanics.}
The Schr\"odinger equation describes the evolution of a quantum particle:
\[
    i\hbar \frac{\partial \psi}{\partial t} = -\frac{\hbar^2}{2m}\Delta \psi + V\psi.
\]

These examples illustrate the richness and universality of PDEs. This course
will provide the mathematical tools needed to understand and solve them.

\section*{References}

This course draws from several standard references:
\begin{itemize}
    \item \textsc{Evans}, L.C., \emph{Partial Differential Equations}, AMS, 2010.
    \item \textsc{Br\'ezis}, H., \emph{Functional Analysis, Sobolev Spaces and
          PDEs}, Springer, 2011.
    \item \textsc{Strauss}, W., \emph{Partial Differential Equations: An
          Introduction}, Wiley, 2008.
    \item \textsc{John}, F., \emph{Partial Differential Equations}, Springer, 1982.
    \item \textsc{Salsa}, S., \emph{Partial Differential Equations in Action},
          Springer, 2016.
    \item \textsc{Renardy}, M. and \textsc{Rogers}, R.C., \emph{An Introduction
          to Partial Differential Equations}, Springer, 2004.
\end{itemize}

\vspace{1cm}
\noindent
Happy reading and good work!

\vspace{0.5cm}
\hfill\emph{The author}
